% Part: normal-modal-logic
% Chapter: completeness
% Section: lindenbaums-lemma

\documentclass[../../../include/open-logic-section]{subfiles}

\begin{document}

\olfileid[pt]{nml}{com}{lin}

\olsection{Lema de Lindenbaum}

O lema de Lindenbaum estabelece que todo conjunto $\Sigma$-consistente
de !!{formula}s está contido em pelo menos um conjunto \emph{completo}
$\Sigma$-consistente. Nossa construção do modelo canônico mostrará que,
para cada conjunto $\Sigma$-consistente completo~$\Delta$, existe um
mundo no modelo canônico no qual são verdadeiras todas e somente as
!!{formula}s em~$\Delta$. Assim, o lema de Lindenbaum garante que todo
conjunto $\Sigma$-consistente é verdadeiro em algum mundo do modelo
canônico.

\begin{thm}[Lema de Lindenbaum]\ollabel{thm:lindenbaum}
  Se $\Gamma$ é $\Sigma$-consistente, então existe um conjunto
  $\Sigma$-consistente completo $\Delta$ que estende~$\Gamma$.
\end{thm}

\begin{proof}
Seja $!A_0$, $!A_1$, \dots{} uma enumeração exaustiva de todas as
fórmulas da linguagem (repetições são permitidas). Por exemplo,
comece enumerando $\Obj p_0$ e, em cada etapa $n \ge 1$, enumere as
finitas fórmulas de comprimento~$n$ que usam somente variáveis entre
$\Obj p_0$, \dots,~$\Obj p_n$. Definimos conjuntos de !!{formula}s
$\Delta_n$ por indução em~$n$ e, em seguida, fazemos
$\Delta = \bigcup_n \Delta_n$. Primeiro, pomos $\Delta_0 = \Gamma$.
Supondo definido $\Delta_n$, definimos $\Delta_{n+1}$ por:
\[
\Delta_{n+1} =
\begin{cases}
  \Delta_n \cup \{!A_n\}, & \text{se $\Delta_n \cup \{ !A_n\}$
    é $\Sigma$-consistente;} \\
  \Delta_n \cup \{ \lnot !A_n\}, & \text{caso contrário.}
\end{cases}
\]
Agora, seja $\Delta = \bigcup_{n=0}^\infty \Delta_n$.

Devemos mostrar que essa definição realmente produz um conjunto
$\Delta$ com as propriedades exigidas, O.s.\@,
$\Gamma \subseteq \Delta$ e $\Delta$ é $\Sigma$-consistente completo.

É evidente que $\Gamma \subseteq \Delta$, pois $\Delta_0 \subseteq
\Delta$ por construção, e $\Delta_0 = \Gamma$. Na verdade, $\Delta_n
\subseteq \Delta$ para todo~$n$, pois $\Delta$ é a união de
todos os~$\Delta_n$. (Como, em cada etapa da construção, acrescentamos
!!a{formula} ao conjunto já construído, $\Delta_n \subseteq
\Delta_{n+1}$; como $\subseteq$ é transitiva, $\Delta_n \subseteq
\Delta_{m}$ sempre que $n \le m$.) Em cada etapa da construção,
acrescentamos ou $!A_n$, ou $\lnot !A_n$, e toda !!{formula} aparece
(pelo menos uma vez) na lista de todos os~$!A_n$. Portanto, para todo
$!A$, ou $!A \in \Delta$, ou $\lnot !A \in \Delta$; logo $\Delta$ é
completo por definição.

Finalmente, devemos mostrar que $\Delta$ é $\Sigma$-consistente. Para
isso, mostramos que (a), se $\Delta$ fosse $\Sigma$-inconsistente,
algum $\Delta_n$ seria $\Sigma$-inconsistente, e (b) todos os $\Delta_n$
são $\Sigma$-consistentes.

Suponha, portanto, que $\Delta$ seja $\Sigma$-inconsistente. Então
$\Delta \Proves[\Sigma] \lfalse$, O.s.\@, existem $!A_1$,
\dots,~$!A_k \in \Delta$ tais que $\Sigma \Proves !A_1 \lif (!A_2
\lif \cdots (!A_k \lif \lfalse)\dots)$. Como
$\Delta = \bigcup_{n=0}^\infty \Delta_n$, cada
$!A_i \in \Delta_{n_i}$ para algum~$n_i$. Seja $n$ o maior desses
índices. Como $n_i \le n$, $\Delta_{n_i} \subseteq \Delta_n$.
Portanto, todos os $!A_i$ estão em algum~$\Delta_n$. Isso significaria
que $\Delta_n \Proves[\Sigma] \lfalse$, O.s.\@, que $\Delta_n$ é
$\Sigma$-inconsistente.

Para mostrar que cada $\Delta_n$ é $\Sigma$-consistente, usamos uma
indução simples em~$n$. Temos $\Delta_0 = \Gamma$, e supusemos que
$\Gamma$ era $\Sigma$-consistente. Portanto, a afirmação vale para
$n = 0$. Suponha agora que ela valha para $n$, O.s.\@, que $\Delta_n$
seja $\Sigma$-consistente. $\Delta_{n+1}$ é
$\Delta_n \cup \{!A_n\}$ se esse conjunto é $\Sigma$-consistente; caso
contrário, é $\Delta_n \cup \{\lnot!A_n\}$. No primeiro caso,
$\Delta_{n+1}$ é claramente $\Sigma$-consistente. Contudo, por
\olref[prf][con]{prop:consistencyfacts}\olref[prf][con]{prop:consistencyfacts-c},
ou $\Delta_n \cup \{!A_n\}$, ou $\Delta_n \cup \{\lnot!A_n\}$ é
consistente, de modo que $\Delta_{n+1}$ também é consistente no outro caso.
\end{proof}

\begin{cor}\ollabel{cor:provability-characterization}
  $\Gamma \Proves[\Sigma] !A$ se e somente se $!A \in \Delta$ para
  cada conjunto $\Sigma$-consistente completo $\Delta$ que estende
  $\Gamma$ (inclusive quando $\Gamma = \emptyset$, caso em que obtemos
  outra caracterização do sistema modal $\Sigma$.)
\end{cor}

\begin{proof}
  Suponha que $\Gamma \Proves[\Sigma] !A$ e seja $\Delta$ qualquer
  conjunto $\Sigma$-consistente completo que estenda $\Gamma$. Se
  $!A \notin \Delta$, então, pela maximalidade, $\lnot!A \in \Delta$;
  assim, $\Delta \Proves[\Sigma] !A$ (pela monotonicidade) e $\Delta
  \Proves[\Sigma] \lnot!A$ (pela reflexividade), e portanto $\Delta$ é
  inconsistente. Reciprocamente, se $\Gamma \Proves/[\Sigma] !A$,
  então $\Gamma \cup \{ \lnot!A\}$ é $\Sigma$-consistente e, pelo
  lema de Lindenbaum, existe um conjunto completo e consistente
  $\Delta$ que estende $\Gamma \cup \{ \lnot!A \}$. Pela consistência,
  $!A \notin \Delta$.
\end{proof}

\end{document}
